%File: paperA_altstyle.tex
%
% !! KEEP THIS FILE PURE ASCII. Overleaf types a file as text or binary when the file
% first enters the project and stores that verdict. A single astral-plane character
% (U+1F534, in a comment) was once enough to register the manuscript as binary and lock
% it against editing, and fixing the bytes afterwards did not clear it. Do not paste
% emoji, typographic dashes, section signs or Greek letters in here, not even inside
% comments. Check before committing with:
%     python paper/asciify_tex.py paper/paperA_altstyle.tex
%     python paper/texlint.py    paper/paperA_altstyle.tex
%
% ALTERNATIVE-STYLE EDITION of Paper A. Same claims, same numbers, same evidence as
% paper/paperA_draft_v2.tex (= main.tex in the PaperA repo); different presentation.
%
% The presentation style is modelled on arXiv:2506.03053v2 (MAEBE): an explicit
% "Research Goals" subsection carrying the questions and claims, a numbered Methods
% section written for replication, a combined "Results and Discussion" section whose
% subsection headings ARE the findings stated as declarative sentences, Related Work
% placed after the results, a single combined "Limitations, Outlook, and Conclusions"
% section, an Impact Statement, and substantive appendices that the main text refers
% forward to.
%
% Constraints held while restyling:
%   - No claim is added, removed, weakened or strengthened. Every number in the v3
%     draft appears here, with the same SRC comment naming its on-disk receipt, and
%     graph_analysis/experiment_paper_claim_audit.py still checks all of them.
%   - The five ML-paper-writing guides the v3 draft was built against still hold:
%     Nanda (claims-first narrative, red-teamed evidence, stated baselines, page-1
%     figure), ICML 2022 Best Practices (claims-to-evidence, reproducibility, compute,
%     licence, metrics, limitations), Foerster (X/Y/Z/1 abstract, contributions as
%     bullets, active voice, citet vs citep, hyperref+cleveref, no filler), Perez
%     (sentence-level prose rules, figure legibility, first-page figure), and
%     Nanda/Perez on empirical alignment research (honest write-up of negative results).
%   - Self-criticism is kept where it bears on a claim (fidelity controls not run,
%     LLM-internal validation, the two reproduction failures, the worked failure case)
%     and dropped where it bears on nothing a reader needs.
%
% Every number carries a source comment: SRC = on-disk receipt in
% graph_analysis/phase2_results/. Markers: [GAP: ...] missing content, [CITE: ...] bib
% entry needed. GAP markers are LaTeX comments only and never render.
%
% CLASS NOTE: deliberately NOT a conference class, for the same reason as the main
% draft: publisher classes carry branding and licence conditions that do not fit an
% arXiv preprint. Stock article, two-column. Switching venue is a preamble-only change.
\documentclass[10pt,twocolumn,letterpaper]{article}

\usepackage[T1]{fontenc}
\usepackage[utf8]{inputenc}
\usepackage{times}
\usepackage[margin=0.75in,columnsep=0.25in]{geometry}
\usepackage{graphicx}
\usepackage{booktabs}
\usepackage{amsmath,amssymb}
\usepackage[hyphens]{url}
\usepackage{natbib}
\usepackage{caption}
\usepackage{abstract}
\usepackage{titlesec}
\usepackage{authblk}
% hyperref late, cleveref after hyperref (Foerster's formatting checklist).
\usepackage[hidelinks,backref=page]{hyperref}
\usepackage[capitalise,noabbrev]{cleveref}

\titleformat{\section}{\normalfont\large\bfseries}{\thesection}{1em}{}
\titleformat{\subsection}{\normalfont\normalsize\bfseries}{\thesubsection}{1em}{}
\titleformat{\subsubsection}{\normalfont\normalsize\itshape}{\thesubsubsection}{1em}{}
\titlespacing*{\section}{0pt}{1.4ex plus .2ex}{0.8ex plus .1ex}
\titlespacing*{\subsection}{0pt}{1.2ex plus .2ex}{0.6ex plus .1ex}
\titlespacing*{\subsubsection}{0pt}{1.0ex plus .2ex}{0.5ex plus .1ex}
\setlength{\parskip}{0pt}
\setlength{\bibsep}{0pt plus 0.3ex}
\renewcommand{\abstractnamefont}{\normalfont\bfseries}
\renewcommand{\abstracttextfont}{\normalfont\small}
\frenchspacing
\pagestyle{plain}

\title{\vspace{-2em}\textbf{Chains, Not Graphs: A Framework for Mechanism-Level\\
Extraction from the AI Safety Literature, and an Evaluation of\\
What the Resulting Graph Does and Does Not Support}\vspace{-0.5em}}

% [GAP: author list + affiliations. Contributors: Martin Leitgab, Gleb Maksimov,
%  Sai (validation), Mike (judge data), Axel (related work), Jeffrey Parks (review).
%  Confirm order + affiliations before submission. Using authblk:
%    \author[1]{Name}  \affil[1]{Institution, \texttt{email}}]
\author{Anonymous Submission}
\date{}

\begin{document}

\twocolumn[
  \begin{@twocolumnfalse}
  \maketitle
  \begin{abstract}
  \noindent
Resources that organise the AI safety literature index it by document metadata or by
topical proximity, so the argument a paper makes internally --- by what mechanism a
proposed intervention is claimed to reduce a specific risk --- is not represented and
cannot be queried across papers. We present an extraction and evaluation framework that
recovers that argument as a typed chain, prompting a reasoning LLM to read 11{,}779
documents of the Alignment Research Dataset and emit, per paper, the path from a named
risk through problem analysis, theoretical insight, design rationale, implementation
mechanism and validation evidence to a proposed intervention, yielding 200{,}525 typed
nodes. Evaluating the result at two levels --- the per-paper chain and the aggregate
graph assembled from all chains --- gives three findings. (1) The chains carry argument
structure rather than a filled-in template: 87.4\% of 2{,}772 quality-controlled
de-duplicated chains contain all five intermediate stages while chain length spans 4--16
nodes, and an LLM-as-judge study over 100 papers with three independent meta-graders
finds omission, not fabrication, to be the dominant error. (2) The chain is a usable
unit of analysis: a risk--intervention query returns one stored chain with its citation
and no clustering step, and the same view shows that one document in six yields a
complete chain and that 2.5\% of extracted interventions are described as deployed.
(3) The aggregate graph does not support the network analysis it invites: merging
near-duplicate nodes manufactures a centrality hub at 90$\times$ the next node that
excluding similarity edges does not repair, and two findings previously reported on this
corpus do not reproduce under controlled re-derivation. We release the pipeline, the
chain dataset, and the audit scripts that regenerate every number reported here.
% All abstract figures audited against raw data 2026-08-11 --
% experiment_paper_claim_audit.json, 42/42 numeric claims PASS.
% The unreproducible 88% figure is deliberately NOT in the abstract.
% See experiment_race_top100_rederive_report.json and paperA_v2_AUDIT.md finding P5.
  \end{abstract}
  \vspace{0.6em}
  \end{@twocolumnfalse}
]

% Figure 1 spans both columns and lands on page 1 (Foerster, Perez).
\begin{figure*}[t]
\centering
\includegraphics[width=\textwidth]{figure1_dataset.png}
\caption{What the extraction produced, and what it does not contain.
\textbf{(Left)} Length distribution of the 2{,}772 de-duplicated quality-controlled
chains. The canonical risk-to-intervention chain has seven nodes (orange); 64.0\% of
chains have exactly that length, 25.8\% are longer because they carry several nodes for
one or more stages, and 10.2\% are shorter and carry genuine gaps. The spread from 4 to
16 nodes is the evidence discussed in \cref{sec:r-template}.
\textbf{(Middle)} LLM-assessed maturity of all 36{,}959 extracted interventions:
roughly six in seven sit at conceptual or early-stage maturity and 2.5\% are described
as deployed. Maturity is an LLM-assigned attribute that our judge protocol does not
score (\cref{sec:m-validation}); we read this panel as coarse composition rather than
as a measured rate.
\textbf{(Right)} Corpus yield: 1{,}868 of the 11{,}779 documents (15.9\%) produce a
complete high-confidence chain, and those documents contribute every chain in the left
panel.}
\label{fig:dataset}
\end{figure*}
% SRC: graph_analysis/experiment_figure1.py, which plots only from committed receipts
%      (experiment_paper_claim_audit.json, experiment_dataset_strength_report.json).
%      No PKL checkpoint required.

%%%%%%%%%%%%%%%%%%%%%%%%%%%%%%%%%%%%%%%%%%%%%%%%%%%%%%%%%%%%%%%%%%%%%%%%
\section{Introduction}
\label{sec:intro}

The AI safety literature grows faster than any individual can track, and the question a
practitioner asks of it is mechanistic: \emph{if I care about risk $R$, what has been
proposed to reduce it, by what mechanism, and how far has that proposal been validated?}
Existing resources cannot answer this question. Curated directories such as the
Alignment Research Dataset (ARD) \citep{kirchner2022ard} and Stampy index papers by
metadata. Embedding-based field maps place papers near topically similar papers. In both
cases the unit of organisation is the document, so the argument inside the document ---
the chain running from a risk to an intervention --- is neither represented nor
queryable, and cannot be compared across papers.

Recovering that argument is harder than the scientific information-extraction tasks that
existing pipelines target. The relation of interest connects concepts that the same
paper introduces rather than two papers or two entities. It spans several reasoning
steps rather than one predicate. And no annotated corpus of such chains exists, so there
is no ground truth to train or evaluate against.

We address this with a framework that has two parts: an extraction stage that prompts a
reasoning LLM (\texttt{o3}) to emit, for each paper, a typed graph running from a
\textsc{risk} endpoint through five intermediate concept subtypes --- \emph{problem
analysis}, \emph{theoretical insight}, \emph{design rationale}, \emph{implementation
mechanism}, \emph{validation evidence} --- to an \textsc{intervention} endpoint, with
per-edge confidence and per-intervention maturity and lifecycle attributes; and an
evaluation stage that measures what the output supports at two levels of aggregation,
the per-paper chain and the graph assembled from all chains. Applied to ARD, the
extraction produces 200{,}525 nodes over 11{,}779 documents.

The two-level evaluation is the part we consider generalisable. A pipeline of this shape
produces both a set of per-document artefacts and, for free, a large graph; the graph
invites clustering, centrality ranking and cross-cluster comparison, and those analyses
are routinely performed. We measure how far each level of aggregation survives scrutiny
on the same extraction, and find that the two levels behave very differently.

\subsection{Research Goals}
\label{sec:goals}

This work answers two questions:

\begin{itemize}
\setlength{\itemsep}{1pt}
\item Does an LLM prompted with a staged schema recover the argument a paper actually
makes, or does it fill in the schema?
\item Which analyses does the resulting artefact support --- at the level of a single
extracted chain, and at the level of the graph assembled from all of them?
\end{itemize}

Our answers form one thesis: \emph{for a corpus of this kind, the per-paper reasoning
chain is a reliable unit of analysis and the aggregate graph is not}. We defend it with
three claims, each of which owns a block of \cref{sec:results}.

\begin{itemize}
\setlength{\itemsep}{2pt}
\item[\textbf{C1}] \textbf{The extracted chains carry argument structure rather than a
filled-in schema.} 87.4\% of 2{,}772 quality-controlled de-duplicated chains contain all
five intermediate reasoning stages while chain length spans 4--16 nodes
(\cref{fig:dataset}, left; \cref{sec:r-stages,sec:r-template}), and an LLM-as-judge
study over 100 papers with three independent meta-graders finds omission, not
fabrication, to be the dominant failure (\cref{sec:r-judge}).
\item[\textbf{C2}] \textbf{Chains are a usable unit of analysis, and reading the corpus
at chain level exposes what it holds.} A risk--intervention query returns one stored
chain with its source citation and no clustering step (\cref{sec:r-retrieval}); the same
view gives a maturity profile and a corpus yield that document-level resources cannot
produce (\cref{sec:r-corpus}).
\item[\textbf{C3}] \textbf{The aggregate graph does not support the network analysis it
invites.} We report four measured artefacts, each with the control that removes it, and
two previously reported findings on this corpus that do not reproduce under controlled
re-derivation (\cref{sec:r-hub}--\cref{sec:r-bydesign}).
\end{itemize}

C3 is a negative result and we report it as a first-class contribution rather than as a
caveat. The quantities from which ecosystem-level claims about this literature are
usually built --- centrality rankings, per-cluster risk-to-intervention ratios,
co-occurrence measured on a centrality-selected head --- are dominated on this corpus by
construction artefacts. We therefore make no ecosystem-level claim ourselves, and
\cref{sec:impact} states what a reader should not do with the released graph.

%%%%%%%%%%%%%%%%%%%%%%%%%%%%%%%%%%%%%%%%%%%%%%%%%%%%%%%%%%%%%%%%%%%%%%%%
\section{Methods}
\label{sec:methods}

This section is written for replication. The extraction pipeline, the path builders, the
audit scripts and the receipt files are released together with the chain dataset
(\cref{sec:m-repro}); every quantity reported in \cref{sec:results} is re-derived from
raw data by a single released script.

\subsection{Corpus Selection}
\label{sec:m-corpus}

We use the Alignment Research Dataset, a curated collection of AI safety publications
spanning arXiv preprints, LessWrong and Alignment Forum posts, EA Forum posts,
organisational reports and transcripts. After filtering to documents that yield a
parseable extraction, the working corpus is 11{,}779 documents.
% SRC: experiment_query_demo_report.json -> coverage_bound.papers_with_risk = 11779

ARD is heterogeneous, and the chain dataset inherits that heterogeneity. Of the 2{,}772
chains defined in \cref{sec:m-reporting}, 1{,}055 come from arXiv, 326 from the
Alignment Forum, 267 from LessWrong, 265 from the EA Forum, 238 from YouTube
transcripts, 120 from MIRI, and the remainder from a long tail of institutional and
personal sites. Seven chains (0.3\%, from three documents) originate from Google Docs.
A ``document'' in ARD is therefore not always a research paper --- the tail includes
reading lists and expository material --- which is why we report source composition
alongside every corpus-level number and why \cref{app:failure} exhibits what the
extraction does with a document of that kind.
% SRC: recomputed 2026-08-11 over paths_hopwise_v4_edge_only_deduped.jsonl,
%      grouping by the source URL host of each chain's risk node.

\subsection{Extraction Schema and Model Configuration}
\label{sec:m-extraction}

Each document is passed in full --- not truncated to its abstract --- to \texttt{o3}
through the OpenAI batch API, with a schema-constrained extraction prompt reproduced in
\cref{app:prompt}. The model traces the paper's own logical chain and emits:

\begin{itemize}
\setlength{\itemsep}{0pt}
\item \textbf{Concept nodes}, each tagged with one \texttt{concept\_category} from
\{\emph{risk}, \emph{problem analysis}, \emph{theoretical insight}, \emph{design
rationale}, \emph{implementation mechanism}, \emph{validation evidence}\}.
\item \textbf{Intervention nodes}, each carrying \texttt{intervention\_lifecycle}
$\in 1..6$ (position in the research-to-deployment cycle) and
\texttt{intervention\_maturity} $\in 1..4$ (1 = conceptual proposal, 2 = early prototype
or experiment, 3 = mature method, 4 = deployed in practice).
\item \textbf{Typed edges} carrying a relation type, a natural-language description of
the claimed relationship, and \texttt{edge\_confidence} $\in 1..5$ reflecting how
explicitly the paper asserts the link.
\end{itemize}

The canonical chain the schema encourages is \emph{risk} $\rightarrow$ \emph{problem
analysis} $\rightarrow$ \emph{theoretical insight} $\rightarrow$ \emph{design
rationale} $\rightarrow$ \emph{implementation mechanism} $\rightarrow$ \emph{validation
evidence} $\rightarrow$ \emph{intervention}. Deviation from it --- omitted stages,
several nodes for one stage --- is permitted by the prompt and is measured in
\cref{sec:r-template} rather than corrected.

Per-paper graphs are embedded (1536-dimensional OpenAI embeddings over node name,
aliases and description) and loaded into a FalkorDB graph database. The result is
200{,}525 nodes: 19{,}096 risks, 36{,}959 interventions, and 144{,}470 intermediate
concept nodes (27{,}748 problem analysis, 26{,}361 theoretical insight, 27{,}543 design
rationale, 34{,}222 implementation mechanism, 28{,}596 validation evidence).
% SRC: experiment_race_prevalence_report.json -> by_node_type_full_graph.*.n
%      (19096+27748+26361+27543+34222+28596+36959 = 200525)

\subsubsection{Extraction is single-source by design}

Every structural edge comes from one paper. The global graph concatenates per-paper
graphs, and the pipeline never asserts a relation between concepts introduced by two
different papers. This is a fidelity choice: an edge in the released graph is a claim
that an identified document made. It has consequences for every downstream graph
statistic, which we make explicit in \cref{sec:m-structural} and measure in
\cref{sec:r-bydesign}.

\subsubsection{Extraction failures and the limits of automated repair}
\label{sec:m-recovery}

Extraction fails outright on a minority of ARD records. Most of those failures are
records with no usable source text --- entries carrying a title and URL but an empty or
\texttt{None} body --- which are unrecoverable in principle. Filtering to failures
repairable in principle (a URL, non-empty source text, and at least one extracted node)
leaves \textbf{441 candidates}. Running the judge of \cref{sec:m-validation} over all
441 in repair mode produced a non-empty graph for \textbf{23 of them (5.2\%)}, and those
recovered graphs are thin: a mean of 2.4 nodes and 0.9 edges, against 16.0 nodes and
10.8 edges for audited successful extractions. An LLM judge audits and repairs an
extraction that already exists effectively, and reconstructs one that failed poorly. We
therefore treat failed extractions as corpus loss rather than as a recoverable pool.
% SRC: experiment_judge_full_report.json ->
%      item3_recovery.population_A_extraction_error_candidates (441 attempts, 23 with a
%      non-empty final_graph = 5.2%, mean 2.35 nodes / 0.87 edges).
% NOTE: the frequently-quoted "~60 recovered of ~400" figure is NOT used here. The 65
%      files in recovered_errors_graph/ are disjoint from these 441 attempts (overlap
%      = 0) and come from a different failure population whose denominator is not in
%      the bundle -- see item3_recovery.PROVENANCE_FINDING.

\subsection{Deduplication and Edge Normalisation}
\label{sec:m-dedup}
% Reused near-verbatim from Gleb_overleaf_paper.txt L143-168.

The raw extraction produces duplicate nodes representing identical concepts with
different phrasings (``US semiconductor export controls'' and ``AI chip export control
regulations'') and inconsistent edge types (\textit{mitigated\_by} versus
\textit{mitigates}). We apply a two-threshold deduplication: candidate pairs require
both cosine similarity $\geq 0.88$ on embeddings and Jaccard overlap $\geq 0.05$ on node
names. We block by node type and category to reduce comparisons, use FAISS
\texttt{IndexIVFFlat} for approximate similarity search, and apply Union-Find to group
transitive duplicates. For merged nodes we select the canonical representative by
longest description and average the embeddings. \Cref{app:dedup} reports the four
threshold configurations tested and the boundary examples that decided the selection. We
prioritise precision over recall, because a false merge conflates distinct concepts
whereas a missed duplicate only introduces redundancy. Edge normalisation unifies
bidirectional pairs and standardises naming variants, yielding 12 canonical edge types.
Deduplication removes 2{,}385 nodes (1.2\%).

\subsubsection{Measured scope of the merge}

The merge is approximate in two respects that we quantify. First, the approximate
nearest-neighbour search under-retrieves: re-running the same threshold specification
with exhaustive search finds 54{,}282 candidate pairs against the 4{,}411 found by the
deployed configuration, a factor of 12.
% SRC: experiment_J6_merge_approx_v2_report.json -> variants.V2_T1_T2_T3.phase2_both_pass
%      = 54282 vs gleb_reference.candidate_pairs = 4411
Second, 1{,}140 nodes across 448 groups remain exact-name duplicates within their
category after the merge.
% SRC: experiment_J6_merge_approx_v2_report.json -> provenance_diagnostic
Both facts matter only where node identity is load-bearing. \Cref{sec:r-hub} shows
exactly where that is; the chain-level results of \cref{sec:r-stages}--\cref{sec:r-corpus}
do not depend on the merge.

\subsection{Structural Diagnostics of the Concatenated Graph}
\label{sec:m-structural}
% Reused from Gleb_overleaf_paper.txt L172-180, relocated from Results to Methods.

The concatenated graph is highly fragmented: 15{,}123 disconnected components, the
largest containing 61 nodes (0.03\%), average degree 2.02, and an average clustering
coefficient of 0.013.
% SRC: re-derived from graph_node_attributes.pkl + graph_edge_data.pkl (EDGE type only)
%      2026-08-11, union-find over 200,525 nodes / 202,149 EDGE edges.
%      NOTE: the frozen Overleaf reported 18,424 / 55 / 2.0 / 0.007 -- measured on the
%      MERGED 200,061-node graph with 197,542 edges, a different substrate. We report
%      our own numbers for the graph this paper describes.
These figures belong in Methods rather than Results because they are the arithmetic
consequence of concatenating per-paper chains without asserting cross-paper edges: a
corpus of $N$ short chains is exactly $N$ components until something links them.

\subsection{Similarity Augmentation and Path Extraction}
\label{sec:m-paths}

Two graph layers are available for analysis, and the distinction governs every claim in
this paper.

\begin{itemize}
\setlength{\itemsep}{1pt}
\item \textbf{EDGE layer (structural).} The extracted relations: single-paper,
LLM-asserted, confidence-scored.
\item \textbf{SIM layer (similarity).} Cosine-similarity edges between nodes of the same
concept category, added to connect content across documents.
\end{itemize}

At $\tau \geq 0.80$ the similarity layer adds 1{,}435{,}806 within-category edges,
reduces the component count from 15{,}123 to 4{,}124, and expands the largest component
from 61 nodes to 152{,}753 (76.2\%). That expansion is the purpose of the layer and also
its hazard: a threshold that turns 0.03\% of the graph into 76.2\% of it in one step is
not a neutral operation. \Cref{app:sensitivity} reports the threshold and hop
sensitivity in full.
% SRC: re-derived 2026-08-11 from graph_edge_data.pkl, SIMILARITY edges with
%      cos_sim >= 0.80 restricted to same concept_category.
%      NOTE: the frozen Overleaf reported 169,083 edges -> 6,522 components ->
%      142,772 (71.4%); that k-NN construction is ~8.5x sparser than the dump this
%      paper uses. Same qualitative conclusion, different counts; we report ours.

For every quantitative result in \cref{sec:r-stages}--\cref{sec:r-corpus} we use the
\textbf{EDGE-only path set} under quality cuts applied at path-construction time:
structural edges only with no similarity hops, edge confidence $\geq 3$,
intervention-endpoint maturity $\geq 3$, and exactly one risk node, which must be the
path root. Enumerating all such paths yields 8{,}954 raw chains. We verify the cuts
empirically rather than trusting the builder configuration: 100\% of enumerated paths
start at a risk node, no path contains a risk after position 0, and every path endpoint
carries maturity 3 or 4.
% SRC: phase2_step4_F2v4_hopwise_falkordb.py (EDGE_CONFIDENCE_MIN=3,
%      INTERVENTION_MATURITY_MIN=3, single-risk); verified in
%      experiment_paper_claim_audit.json.

\subsection{Chain De-duplication and the Reporting Unit}
\label{sec:m-reporting}

Raw path enumeration over-counts prolific papers: one paper contributes 700 of the
8{,}954 raw chains (7.8\%), two others contribute 200 and 184, and 169 papers contribute
ten or more. Per-chain statistics on that set would weight the corpus by how branched
each paper's extracted graph happens to be. We therefore collapse near-duplicate chains
within each source paper: two chains are near-duplicates when the smaller one's node set
is $\geq 70\%$ contained in the larger's, and the shorter is dropped. This yields the
\textbf{2{,}772-chain de-duplicated set} (at most 7 chains from any one paper), which is
the reporting unit for every number below. Both sets cover the same \textbf{1{,}868
source papers}, and we report the raw set alongside as a sensitivity check.
% SRC: phase1_dedup_paths.py (CONTAINMENT_THRESHOLD = 0.70 -- the module docstring says
%      80%, the code uses 70%; the code is what ran).
%      Counts and skew: experiment_paper_claim_audit.json -> path_set_comparison.

The path set is free of similarity-threshold sensitivity and of merge artefacts by
construction. It is the substrate on which we make numerical claims.

\subsection{LLM-as-Judge Validation Protocol}
\label{sec:m-validation}

We assess extraction quality with an LLM-as-a-judge protocol \citep{zheng2023judging}
and a judge-of-judge calibration stage.

\paragraph{Judge configuration.} A judge model from a different provider than the
extractor (\texttt{claude-sonnet-4-5}, batch API) receives the source text together with
the extracted graph and returns a structured audit: schema conformance, referential
integrity, orphan and duplicate detection, rationale consistency, coverage of the
source's claims, and a set of proposed repairs. We run it on a random sample of
\textbf{100 successfully extracted papers} spanning all nine ARD source types; the
per-source-type breakdown is in \cref{app:judge}.

\paragraph{Meta-grader configuration.} Three meta-graders
(\texttt{claude-opus-4-5}, \texttt{gemini-3-pro}, \texttt{gpt-5.1}) score each extraction
on a 0--100 holistic rubric before and after the judge's proposed repairs. The graders
receive the source text, the extraction, and the judge report including its
\emph{proposed} fixes. The post-repair score therefore assesses the proposed repairs as
described rather than measuring a rebuilt graph. Two consequences follow, and both
matter for reading \cref{sec:r-judge}: no repaired graph was re-extracted and re-scored,
and defects in the fix-application code cannot have contaminated these scores, because
no applied graph reached a grader.
% SRC: prompts/judge_and_extraction_evaluation_prompt.md ("give a score between 0-100
%      before and after the judge's proposed fixes") + combine_judge_and_extraction_
%      with_original_text.py, which bundles Original Text + Extraction Output + Judge
%      Output for the meta-grader. Verified 2026-08-11.

\paragraph{Scope of the protocol.} The judge scores node and edge faithfulness,
referential structure and edge type. It does \emph{not} score the
\texttt{intervention\_maturity} assignment, so maturity is labelled LLM-assessed
wherever it appears and is used descriptively. No human adjudicated any judge or
meta-grader verdict, so the validation chain is LLM-internal end to end; a
human-anchored spot-check and a manual error taxonomy are the outstanding items
(\cref{sec:limitations}).
% [GAP: human-anchored spot-check + manual 50-instance error taxonomy. Owner: Sai,
%  ticket #150. When they land, they attach to sec:r-judge and to sec:limitations.]

\subsection{Reproducibility, Release and Compute}
\label{sec:m-repro}

Every quantitative claim in this paper carries a source comment in the manuscript naming
the receipt file it came from, and one released script re-derives all of them from the
raw path files and the node attributes, then compares each against the value printed
here. On the current draft it reports 42 of 42 numeric claims passing. A
claim-to-script-to-receipt map ships with the code.

The extraction pipeline, the path builders, the audit scripts and the receipt files are
released with the chain dataset under a permissive licence. Two large intermediate
checkpoints (node attributes and edge data, together $\sim$3.2\,GB) are rebuildable from
the released graph dump with a released script; because the receipts are committed,
every number here is checkable without them. \Cref{fig:dataset} plots only from
committed receipts for the same reason.

\paragraph{Compute.} Extraction used the OpenAI batch API with \texttt{o3} over 11{,}779
documents. The judge and meta-grader runs used the Anthropic, OpenAI and Google batch
APIs at the model versions named above; model identifiers are pinned rather than
floating, so a re-run reproduces the same model generation. All graph analysis runs on a
single workstation CPU without a GPU, with path enumeration against the graph database
as the heaviest step.
% [GAP: exact wall-clock and token totals for the extraction run, if recoverable from
%  the batch-API logs. Not blocking.]

%%%%%%%%%%%%%%%%%%%%%%%%%%%%%%%%%%%%%%%%%%%%%%%%%%%%%%%%%%%%%%%%%%%%%%%%
\section{Results and Discussion}
\label{sec:results}

Results are organised by level of aggregation. \Cref{sec:r-stages}--\cref{sec:r-judge}
report properties of individual extracted chains (C1);
\cref{sec:r-retrieval}--\cref{sec:r-corpus} report what the chain supports as a unit of
analysis (C2); \cref{sec:r-hub}--\cref{sec:r-bydesign} report what the aggregate graph
does not support (C3). Every number in the first two blocks is a structural property of
the EDGE-only quality-cut path set (\cref{sec:m-paths}) or a direct count over extracted
attributes: no clustering, centrality or similarity threshold enters any of them.

\paragraph{Provenance of the C3 analyses.} The measurements in
\cref{sec:r-hub}--\cref{sec:r-bydesign} are \emph{post-hoc}. An earlier analysis pass
over this corpus produced a set of graph-analytic findings; we re-derived them under
controls, and two did not survive. We report the re-derivations, and we flag the
post-hoc status because a hypothesis formed after seeing the data deserves less weight
than one registered before it.

\subsection{87.4\% of extracted chains carry all five reasoning stages}
\label{sec:r-stages}

Across the 2{,}772 de-duplicated chains, \textbf{87.4\% (2{,}423) contain all five
intermediate reasoning stages}. Stage-wise presence is 99.0\% for problem analysis,
93.9\% for theoretical insight, 96.1\% for design rationale, 96.0\% for implementation
mechanism, and 98.0\% for validation evidence. On the raw 8{,}954-chain set the same
figure is 89.0\%, so the result does not depend on the de-duplication choice.
% SRC: experiment_paper_claim_audit.json -> path_set_comparison.deduped
%      (all5 2423/2772 = 87.4%; subtype_presence_pct) and .raw (89.0%)

A large majority of extracted chains therefore do more than connect a risk to an
intervention: they carry the paper's argument for why the intervention should work,
stage by stage.

This percentage is a share \emph{of complete chains}, not of the corpus. Only 1{,}868 of
the 11{,}779 documents (15.9\%) yield a complete high-confidence chain at all; the
remainder contribute nodes and partial structure. Both numbers matter, and
\cref{sec:r-corpus} reports the corpus-level view.

\subsection{Chain length spans 4 to 16 nodes, so the schema is not a filled-in template}
\label{sec:r-template}

The obvious objection to \cref{sec:r-stages} is that a schema suggesting a seven-node
chain will produce seven-node chains regardless of the source. \Cref{fig:dataset} (left)
refutes the strong form of that objection: chains span 4--16 nodes; the mode is 7
(1{,}775 chains, 64.0\%); 25.8\% run \emph{longer} than the canonical seven, carrying
several nodes for one or more stages, that is, richer argumentation; and 10.2\% run
shorter, carrying genuine gaps.
% SRC: experiment_paper_claim_audit.json -> path_set_comparison.deduped.len_hist
%      (=7: 1775 = 64.0%; >7: 715 = 25.8%; <7: 282 = 10.2%)

\subsubsection{What the length spread does and does not settle}

A spread is consistent with extraction following each paper's argument, and it is also
consistent with a schema that bends a little under pressure. Settling the question
requires separating two failure modes that are easy to conflate. A \emph{scope} failure
is a chain emitted from a document that is not about AI safety. A \emph{fidelity}
failure is a chain whose structure the source document does not contain. Only the second
bears on C1.

An out-of-domain control does not test fidelity. Any literature that names a problem,
analyses it and proposes a remedy --- materials science, epidemiology, software
engineering --- contains genuine risk-to-intervention arguments, so a high emission rate
on non-safety documents would show that the schema is domain-general, which is a scope
property and arguably a feature. Two controls do test fidelity, and both hold the domain
fixed:

\begin{itemize}
\setlength{\itemsep}{1pt}
\item \textbf{Schema ablation.} Re-extract a sample with a prompt that asks for the
chain from problem to proposed solution \emph{without naming the five stages}, then
measure how often the emergent chain maps onto them. Structure that survives un-prompted
is read from the document rather than imposed on it. This control doubles as the
ablation that \cref{sec:limitations} identifies as missing.
\item \textbf{Degraded-source control.} Re-extract from versions of the same safety
documents whose argument is destroyed but whose vocabulary is retained ---
sentence-shuffled text, abstract only, reference list only. Complete five-stage chains
arising at a similar rate would show the model confabulating structure from topical
vocabulary.
\end{itemize}

We ran neither, and we identify them as the highest-value additions to this work.
\Cref{app:failure} exhibits the fidelity failure we can demonstrate: a chain that passes
every structural filter, extracted from an expository article proving that there are
infinitely many primes, in which the model invents a safety framing the source never
asserts. It is an existence proof, not a rate.

\subsubsection{What C1 does not claim}

C1 concerns fidelity to what a paper argued. It says nothing about whether that argument
is \emph{correct}, whether the intervention works, or whether the paper is good safety
research. A faithful extraction from a weak paper is a successful extraction.

\subsection{The judge study finds omission, not fabrication, to be the dominant error}
\label{sec:r-judge}

\paragraph{Structural error classes.} Across the 100 audited papers the judge found
\textbf{108 referential-integrity findings}, of which 68 are blocker-severity, the
dominant pattern being an edge whose source or target names a node absent from the node
list; \textbf{42 orphan nodes}, present in 40\% of papers; and \textbf{56 within-paper
duplicate node pairs}, in 55\% of papers. The mean audited graph has 16.0 nodes and 10.8
edges.
% SRC: experiment_judge_item2_report.json -> findings_total,
%      findings_papers_with_at_least_one, per_paper_means

\paragraph{Coverage.} Asked to identify content in the source that the extraction
missed, the judge proposed \textbf{new nodes for only 7 of 100 papers} (9 nodes in
total). Its coverage list flags a mean of 7.8 relationships per paper that it judges the
source to support and the extraction not to contain --- an upper-bound opinion signal
rather than a measured omission rate, since the judge has no schema slot in which to
propose those edges and no human adjudicated them.
% SRC: experiment_judge_item2_report.json -> proposed_fix_totals.add_nodes = 9;
%      findings_papers_with_at_least_one.add_nodes = 7;
%      per_paper_means.missing_edges_flagged

\paragraph{Direction of the judge effect is robust; magnitude is not.}
\Cref{tab:graders} reports the meta-grader scores. All three graders agree on direction:
every grader records an improvement on every paper it scored, and on the 13 papers all
three scored, agreement that the judge improved the extraction is unanimous. They
disagree sharply on magnitude, with mean gains from $+3.5$ to $+33.1$ points.
Inter-grader agreement on those 13 papers is moderate before the repairs (Fleiss'
$\kappa = 0.54$ on four ordinal score bands) and collapses afterwards
($\kappa = 0.09$): Opus and GPT-5.1 remain rank-correlated on the post-repair scores
(Spearman $\rho = 0.50$) while Gemini's saturated scores correlate negatively with both
($\rho = -0.63$ and $-0.64$). We report the \emph{direction} of the judge's effect as a
robust, unanimous finding and treat its \emph{magnitude} as unresolved.
% SRC: experiment_judge_full_report.json -> item2_meta_graders, item2_grader_agreement
%      (reproduces extraction_validator/results.md exactly)

\begin{table}[t]
\centering
\small
\begin{tabular}{lrrrr}
\toprule
\textbf{Meta-grader} & \textbf{$n$} & \textbf{pre} & \textbf{post} & \textbf{$\Delta$} \\
\midrule
Claude Opus 4.5 & 95 & 69.7$\pm$18.0 & 73.1$\pm$18.2 & $+3.5$ \\
GPT-5.1         & 95 & 67.0$\pm$16.2 & 80.2$\pm$10.2 & $+13.2$ \\
Gemini 3 Pro    & 13 & 62.7$\pm$22.6 & 95.8$\pm$1.8  & $+33.1$ \\
\bottomrule
\end{tabular}
\caption{Meta-grader rubric scores before and after the judge's proposed repairs, mean
$\pm$ standard deviation. The graders scored unequal, partly disjoint subsets of the 100
papers --- only 13 papers carry all three scores --- because the rubric ran across
several prompt iterations and only some iterations emit a paired pre/post score. Rows
are therefore not directly comparable, and Gemini's $n{=}13$ post-repair mean of 95.8
with a standard deviation of 1.8 indicates saturation at the top of the scale.}
\label{tab:graders}
\end{table}

\paragraph{Error profile.} The Opus meta-grader recorded 216 missed concepts across the
95 papers it scored, at least one in every paper, against 20 instances of fabricated
content in 16\% of papers and 15 category errors in 33\%.
% SRC: experiment_judge_full_report.json -> item3_error_taxonomy_auto
Under this grader the dominant failure mode is \emph{omission, not hallucination}. That
supports reading an extracted chain as a partial record of what a paper argued, and does
not support reading the absence of a stage as evidence that the paper lacked one. These
counts are un-adjudicated LLM output and we report them as a profile rather than as a
measured error rate.

\paragraph{One measurement artefact to separate before quoting any error rate.} The
judge grades against a schema expecting rationale fields inline on nodes and edges,
whereas the pipeline stores rationale as separate nodes linked by a dedicated relation.
That mismatch accounts for \textbf{184 of the 218 blocker-severity schema flags
(84\%)}; excluding it leaves 34 substantive blocker-severity schema findings across 100
papers. A headline error rate computed from raw judge output without this split measures
a schema version difference rather than extraction quality.
% SRC: experiment_judge_item2_report.json ->
%      schema_check_severity_rationale_field_mismatch vs schema_check_severity_substantive

\subsection{A risk--intervention query returns one stored chain with its citation}
\label{sec:r-retrieval}

The dataset supports a query that document-level resources cannot serve: \emph{by what
mechanism does intervention $I$ reduce risk $R$?} Because the mechanism is stored as an
explicit path, answering it is a single graph lookup returning the chain and its source
citation, with no clustering, no ranking and no embedding search over abstracts.
\Cref{tab:query} shows two chains retrieved this way, one governance and one technical,
both from arXiv sources so that the extraction input is structured and the citation is
unambiguous.
% SRC: experiment_query_demo_report.json -> query_demo.*

\begin{table*}[t]
\centering
\small
\begin{tabular}{p{0.09\textwidth}p{0.86\textwidth}}
\toprule
\textbf{Stage} & \textbf{Query 1 --- governance.} Source: \citet{shavit2023chinchilla} \\
\midrule
risk & Covert development of dangerous ML models using large-scale training \\
pa & Tampering with logging mechanisms via supply chain attacks \\
ti & Open-source Hardware Roots-of-Trust enhance transparency and tamper evidence \\
dr & Embed open-source HRoT like Caliptra in ML chips \\
im & Signed firmware enforced by HRoT with anti-rollback protections \\
va & Industry adoption of Caliptra HRoT across major vendors \\
\textbf{intv} & Adopt open-source HRoT in data-center ML chips to enforce secure logging \textbf{(maturity 3)} \\
\midrule
\multicolumn{2}{c}{} \\
\midrule
\textbf{Stage} & \textbf{Query 2 --- technical alignment.} Source: \citet{ngo2022alignment} \\
\midrule
risk & Catastrophic disempowerment of humanity by misaligned AGI \\
pa & Misaligned goal formation during AGI training \\
ti & Reward misspecification in real-world RL tasks \\
dr & Automated scalable supervision to reduce reward misspecification \\
im & Reinforcement learning from human feedback scaling \\
va & Instruction-following improvement in InstructGPT via RLHF \\
\textbf{intv} & Scale RLHF fine-tuning using AI-generated feedback signals \textbf{(maturity 3)} \\
\bottomrule
\end{tabular}
\caption{Two mechanism-level retrievals, one governance and one technical, both from
arXiv sources. Each block is a single extracted chain returned verbatim by a
risk--intervention query; stage abbreviations follow the logical chain order
(pa = problem analysis, ti = theoretical insight, dr = design rationale,
im = implementation mechanism, va = validation evidence). Neither answer is synthesised:
each is one stored path with its source document attached.}
\label{tab:query}
\end{table*}
% SRC: Query 1 = node ids [1067,1086,1087,1088,1089,1090,1091], arxiv.org/abs/2303.11341
%      Query 2 = node ids [415,416,417,426,431,437,442],       arxiv.org/abs/2209.00626
%      Both verified against paths_hopwise_v4_edge_only_deduped.jsonl 2026-08-11.
%      Titles/authors verified against arXiv the same day. See paperA_v2_AUDIT.md P3.

Three properties of these results matter. They are \emph{single-hit}: the answer is one
stored path, not a synthesis over retrieved passages. They are \emph{attributable}:
every node carries its source document. And they are \emph{comparable}: because all
chains share the same stage vocabulary, chains from different papers addressing the same
risk sit side by side, which is the precondition for the cross-paper mechanism analysis
identified in \cref{sec:outlook}.

\subsection{One document in six yields a complete chain, and 2.5\% of interventions are described as deployed}
\label{sec:r-corpus}

The same unit supports descriptives that no metadata directory can produce, because they
are properties of the argument inside each paper.

\paragraph{Intervention maturity.} Of the 36{,}959 extracted interventions, the
LLM-assessed maturity distribution is \textbf{27.6\% conceptual} (10{,}189),
\textbf{57.3\% early-stage} (21{,}180), \textbf{12.7\% mature} (4{,}681) and
\textbf{2.5\% deployed} (909); see \cref{fig:dataset} (middle). Roughly six in seven
proposed interventions in this corpus sit at conceptual or early-stage maturity, and
fewer than one in thirty is described as deployed. Maturity is an LLM assignment that no
human adjudicated and that the judge protocol does not score (\cref{sec:m-validation}),
so we read it as coarse composition. The skew is a four-bucket call with a very large
margin and would survive a substantial per-item error rate; the individual bucket
boundaries should not be read as precise.
% SRC: experiment_dataset_strength_report.json -> interventions.maturity
% NOTE 2026-08-11: a scope-composition paragraph (share of chains whose risk is a
%      capability gap rather than a human harm) was drafted here and WITHDRAWN. It
%      depended on phase2_routing_assignments.jsonl, which is not committed and is not
%      re-derivable by a reader at reasonable cost, and it belongs to a separate line of
%      work. Do not reinstate it unless that file ships with the release and a committed
%      script reproduces it. See paperA_v3_RESTRUCTURE.md Sec.2.2.

\paragraph{End-to-end coverage.} All 11{,}779 corpus documents yield at least one risk
node and at least one intervention node --- the two sets are identical --- which is a
property of the extraction schema, not a measurement of the literature. The informative
number is how many documents yield a \emph{complete, high-confidence,
mature-intervention chain end to end}: \textbf{1{,}868 documents} (15.9\%), contributing
\textbf{2{,}643 distinct risk-to-intervention node pairs} after de-duplication (3{,}222
before); see \cref{fig:dataset} (right). The ratio between the two figures is the honest
statement of yield: roughly one document in six argues a complete mechanism from a named
risk to an intervention that is at least a mature method, under our quality cuts.
% SRC: experiment_paper_claim_audit.json -> n_source_papers = 1868 recomputed from the
%      raw path file (identical on the de-duplicated set); distinct_ri_node_pairs
%      deduped 2643 / raw 3222.
%      NOTE: experiment_query_demo_report.json carries "papers_with_complete_strict_
%      chain": 1868 but it is a HARDCODED literal in experiment_query_demo.py:125, not
%      computed there. Cite the audit receipt.

\subsection{Merging near-duplicate nodes manufactures a centrality super-hub}
\label{sec:r-hub}

A graph of this shape invites standard network analysis: cluster the nodes, rank by
centrality, compare risk and intervention densities per cluster, look for co-occurring
framings. This block reports what each of those quantities measures once construction
artefacts are controlled for. Each result is stated as a measurement on this corpus
together with the control that removes it.

The corpus contains thousands of near-identical existential-risk concept nodes
(``Existential catastrophe from misaligned advanced AI'' and close paraphrases)
contributed by different papers. Merging them is the natural cleanup step, and it is
what breaks centrality, because the canonical node inherits every member's structural
edges. In our measurement the largest merged risk node absorbed \textbf{4{,}066 members
and 7{,}777 structural edges}.
% SRC: experiment_merge_vs_simexcl_ec_report.json ->
%      biggest_merged_node_members / _edge_degree

\Cref{tab:ec} reports eigenvector centrality under the four combinations of (merged /
un-merged) $\times$ (similarity edges included / risk--risk similarity excluded).

\begin{table}[t]
\centering
\small
\begin{tabular}{lrrr}
\toprule
\textbf{Condition} & \textbf{EC$_1$} & \textbf{EC$_1$/EC$_2$} & \textbf{xrisk} \\
                   &                 &                        & \textbf{in top-10} \\
\midrule
un-merged, full SIM             & 0.0328 & 1.0$\times$ & 10/10 \\
un-merged, risk--risk SIM excl. & 0.0011 & 1.0$\times$ & 8/10 \\
merged, full SIM                & 0.5069 & 33.6$\times$ & 1/10 \\
merged, risk--risk SIM excl.    & 0.3337 & \textbf{90.1$\times$} & 2/10 \\
\bottomrule
\end{tabular}
\caption{Eigenvector centrality of the top node under four graph-construction choices.
Merging near-duplicate risk nodes produces a single node with 90$\times$ the centrality
of the next node, and excluding similarity edges does not repair it. On the un-merged
graph the top nodes fall within 1\% of each other, so no ranking exists to report.}
\label{tab:ec}
\end{table}
% SRC: experiment_merge_vs_simexcl_ec_report.json -> conditions[0..3].top10[0].ec and
%      ratio to top10[1].ec: 0.032753/0.032686=1.00 ; 0.0010682/0.0010393=1.03 ;
%      0.50685/0.015084=33.6 ; 0.33369/0.0037034=90.1

Two consequences follow, and their order matters. \textbf{(a) Un-merging is primary.} On
the merged graph the super-hub retains a 90$\times$ lead even after deleting all
risk--risk similarity edges, because the merge baked the concentration into the
structural topology. \textbf{(b) Excluding risk--risk similarity is still necessary},
but as the \emph{second} step: on the un-merged graph it flattens the top centrality
value by a factor of 31 (0.0328 to 0.0011) by breaking the near-duplicate similarity
clique. Applied alone, either control leaves the artefact intact.

The substantive point for readers of this literature is that ``the most central risks
are existential'' is a statement about graph construction rather than about the field.

\paragraph{Control.} Do not merge near-duplicates before computing centrality. Exclude
within-category similarity edges from all centrality measures. Count canonical concepts
or source papers, never raw nodes, for any co-occurrence statistic.

\subsection{Approximate duplicate search under-removes by a factor of twelve}
\label{sec:r-recall}

Deduplication quality is usually reported as a removal rate. That number is uninformative
on its own: re-running our own threshold specification with exhaustive rather than
approximate search returns \textbf{54{,}282 candidate pairs against 4{,}411}, a factor of
12, and the gap concentrates almost entirely in the risk block --- precisely where the
centrality artefact of \cref{sec:r-hub} lives.
% SRC: experiment_J6_merge_approx_v2_report.json -> variants.V2_T1_T2_T3 vs gleb_reference

\paragraph{Control.} Report recall of the candidate-generation stage, not only the
removal rate, and verify recall on the node block to which downstream metrics are most
sensitive.

\subsection{Node-level clustering groups topics, not mechanisms}
\label{sec:r-clustering}

Direct $k$-means on the 1536-dimensional embeddings fails outright (silhouette
$\approx 0.02$). UMAP reduction to 150 dimensions followed by $k$-means at $k{=}40$
reaches silhouette 0.298 --- modest against conventional interpretation thresholds
\citep{kaufman1990finding}; \cref{app:clustering} gives the full comparison --- and the
resulting clusters are recognisable field topics (\cref{app:clusters}). But cluster
co-membership at $\tau{=}0.80$ tracks vocabulary rather than mechanism: the same
threshold that produces navigable topics also grows the largest connected component from
61 nodes to 152{,}753. We attempted to derive a clean, non-overlapping risk and
intervention taxonomy this way, tightening quality cuts step by step --- maturity, edge
confidence, EDGE-only paths, clustering each node type on its own axis --- and
node-level semantic clustering remained insufficient to carry mechanism-level
conclusions.

\paragraph{Control.} Treat node-level clusters as a browsing index, not as an
inferential unit. Cluster risk and intervention nodes on separate axes. Do not compute
risk-to-intervention ratios per cluster, because numerator and denominator are populated
by different extraction stages and the ratio inherits the merge artefact of
\cref{sec:r-hub}.

\subsection{Co-occurrence measured on a centrality-selected head does not reproduce}
\label{sec:r-selection}

Selecting the top-100 risks by eigenvector centrality and measuring which framings appear
there is a natural exploratory move, and in a previously reported analysis of this corpus
it produced a striking result: 88\% of the single-path risks in that head appeared to be
framed through competitive race dynamics, with an apparent 44-fold gradient from the
centrality head to the full risk population.

Neither figure survives re-derivation. Re-running the measurement under the same
definitions --- importance by eigenvector centrality, path diversity by distinct
first-hop structural problem-analysis neighbours, race framing by the same keyword rule
--- reproduces the \emph{structural} part (38 of the top-100 risks are single-path,
against 41 reported) but not the framing part: 2.6\% of those single-path risks are
race-framed, and the head-to-population gradient is 2.0-fold, not 44-fold. The keyword
classifier itself is sound --- \cref{app:race} reports 92\% precision on a stratified
manual sample --- so the failure lies in the selection step rather than in the
measurement instrument.
% SRC: experiment_race_top100_rederive_report.json -> condition_A (merged + full
%      within-category SIM, the closest reconstruction of the original setup). The
%      similarity layer in the graph dump we use is ~8.5x denser than the one the
%      original analysis ran on, so the centrality ranking -- and hence the selected
%      head -- is not identical. We report the failure to reproduce rather than the
%      original figure.

Corpus-wide prevalence of the framing is \textbf{2.21\% of problem-analysis nodes}
(614/27{,}748) and 1.65\% of risk nodes; only 0.50\% of EDGE-only chains contain a
race-framed problem-analysis node, and the re-derived head rates (1.3--6.3\% across
centrality tiers) sit within a small multiple of that base.
% SRC: experiment_race_prevalence_report.json + experiment_race_top100_rederive_report.json

A related quantity behaves the same way: deleting race-framed problem-analysis nodes
appeared to disconnect 51 of the top-100 risks from all interventions. Recomputed on the
corrected merge, the same deletion isolates \textbf{2 of 100}.
% SRC: experiment_J6_merge_approx_v2_report.json -> j6_on_V2.isolated_after_race_removal
%      = 2 vs gleb_reference.j6_isolated = 51

\subsubsection{The lesson is stronger, not weaker, for the reproduction failing}

A selection-conditioned statistic on this kind of graph is not merely inflated, it is
\emph{unstable}: it depends on the similarity-layer density and on the merge decisions
that set the ranking, and those are pipeline parameters rather than properties of the
literature. No number computed this way should be reported without the selection rule,
the graph construction, and a full-population baseline alongside it.

\subsubsection{The underlying association is real but modest}

\Cref{tab:race} reports the race-by-existential-risk co-occurrence odds ratio under three
units of analysis. The association survives collapsing near-duplicates, which rules out
pure pseudo-replication, but at 2.7--5.2$\times$ on a $\sim$2\% base rate it supports no
claim about the discourse structure of the field.

\begin{table}[t]
\centering
\small
\begin{tabular}{lrrr}
\toprule
\textbf{Unit} & \textbf{OR} & \textbf{95\% CI} & \textbf{$P$(race)} \\
\midrule
node               & 3.82 & 3.24--4.50 & 7.6\% vs 2.1\% \\
merged node        & 2.73 & 2.06--3.64 & 5.5\% vs 2.1\% \\
source paper       & 5.17 & 4.32--6.20 & 10.7\% vs 2.3\% \\
\bottomrule
\end{tabular}
\caption{Race-framing by existential-risk co-occurrence across units of analysis (all
$p < 10^{-12}$). The $P$(race) column gives the rate within existential-risk versus
non-existential-risk nodes. The association is directionally robust and quantitatively
minor.}
\label{tab:race}
\end{table}
% SRC: experiment_race_dedup_robustness_report.json ->
%      A_node_level / B_merged_node_level / C_paper_level

\paragraph{Control.} Report prevalence on the full population before reporting any
pattern in a selected head. State the selection rule as part of the result. Choose the
unit of analysis --- node, merged node, paper --- explicitly, and show that the effect
survives all three.

\subsection{Single-source extraction makes several structural statistics true by design}
\label{sec:r-bydesign}

Three quantities that look like discoveries are consequences of the pipeline.
\emph{Path diversity}: because each paper contributes one chain, a risk discussed in one
paper necessarily has one route to an intervention, which is extraction structure rather
than fragility in the literature. \emph{Path completion}: under risk-rooted
breadth-first search over a similarity-augmented graph most branches terminate at
intermediate stages, so a low completion percentage measures the traversal rather than
the corpus. \emph{Siloing}: cross-paper links exist only in the similarity layer, so the
structural layer is siloed by construction while the similarity layer's connectivity is a
function of $\tau$, and no layer measures real siloing.

\paragraph{Control.} Before reporting any structural statistic, ask what value it would
take on a corpus of independent single-paper chains. If the answer is the value you
observed, it is a pipeline property.

%%%%%%%%%%%%%%%%%%%%%%%%%%%%%%%%%%%%%%%%%%%%%%%%%%%%%%%%%%%%%%%%%%%%%%%%
\section{Related Work}
\label{sec:related}
% [GAP: Axel started this section. The five strands below are the agreed scope and all
%  citations resolve; a final read-through for flow is the remaining work.]

\paragraph{Knowledge graphs from scientific literature.} Prior work extracts entities and
relations from scientific text, typically targeting citation, method--task--dataset, or
claim structures. \citet{luan2018scierc} is representative, jointly identifying entities,
relations and coreference to build a scientific knowledge graph, and the task has a long
methodological history \citep{bach2007relation}. Our schema differs in that the relation
of interest is an \emph{argument} internal to one paper --- risk to intervention through
explicit reasoning stages --- rather than a bibliographic or method-level relation.

\paragraph{Agentic and self-organising knowledge graphs.} A parallel line of work has
models build and extend knowledge graphs autonomously. \citet{buehler2025agentic} grows a
scale-free knowledge network through iterative model feedback, and
\citet{rasmussen2025zep} uses a temporal knowledge graph as agent memory. Those systems
optimise a graph for an agent's downstream reasoning; we optimise for fidelity to what a
specific paper argued, which is why our edges are confidence-scored assertions traceable
to one source document rather than model-generated links.

\paragraph{Graph-based retrieval-augmented generation.} GraphRAG systems build a graph
over a corpus so that a language model retrieves structured context rather than isolated
passages \citep{peng2024graphrag,sarmah2024hybridrag}; the design space has been mapped
modularly \citep{cao2024legographrag} and pushed towards learning graph construction
against downstream task reward \citep{tsang2025autographr1}. \citet{pan2024unifying}
survey the broader programme of joining language models to knowledge graphs, and
\citet{wu2025medgraphrag} build a medical GraphRAG whose selling point is credible
sourcing. In all of that work the graph is \emph{instrumental}: it exists to improve an
answer, and its edges are whatever the indexing step produced. Our graph is the object of
study, each edge is a confidence-scored assertion that one identified paper made, and a
query returns that paper's own argument with provenance attached rather than a
synthesised answer. A GraphRAG layer over this dataset is an obvious application we do
not attempt here.

\paragraph{AI safety literature resources.}
% [CITE: AI Safety Atlas -- no canonical reference resolved; cite URL inline if used.]
The Alignment Research Dataset \citep{kirchner2022ard}, maintained by the Stampy /
AI Safety Info project, and hand-curated topic maps such as the AI Safety Atlas organise
the field at the level of documents or topics. The MIT AI Risk Repository
\citep{slattery2026airisk} provides a curated risk vocabulary --- 1{,}725 risks distilled
from 74 frameworks into causal and domain taxonomies. These complement rather than
compete with a mechanism-level extraction: a risk taxonomy says what can go wrong, a
directory says who wrote about it, and neither records the argument connecting a risk to
a proposed remedy. Anchoring our extracted risks to an external vocabulary of this kind
is the natural next step (\cref{sec:outlook}).

\paragraph{Causal and mechanism extraction.}
% [CITE: PICO extraction in evidence-based medicine; causal claim extraction.]
Evidence-based medicine has a mature tradition of extracting
population--intervention--comparison--outcome structure from trials. Our task is the
analogous reduction for a field whose evidence is largely theoretical, which is why
validation evidence appears as a reasoning stage rather than as an outcome measurement.

\paragraph{LLM extraction validation.} We adopt the LLM-as-a-judge paradigm
\citep{zheng2023judging} with an explicit judge-of-judge calibration step, in which
several independent stronger models score the judge's proposals.
\Cref{sec:limitations} states its scope limits.

%%%%%%%%%%%%%%%%%%%%%%%%%%%%%%%%%%%%%%%%%%%%%%%%%%%%%%%%%%%%%%%%%%%%%%%%
\section{Limitations, Outlook, and Conclusions}
\label{sec:lim-out-conc}

\subsection{Limitations}
\label{sec:limitations}

\paragraph{No comparative baseline.} We claim that this extraction recovers argument
structure, not that it outperforms an alternative extractor. We compare against no
simpler pipeline --- flat triple extraction, abstract-only extraction, a non-reasoning
model --- and we ran no ablation over schema components. A reader should therefore read
C1 as evidence that the output has the claimed structure, not as evidence that the design
choices are necessary for it. The schema-ablation and degraded-source controls of
\cref{sec:r-template} are the highest-value missing experiments, and the first doubles as
the missing ablation.

\paragraph{Validation is LLM-internal.} An LLM judges faithfulness with judge-of-judge
calibration; domain experts do not, at any scale. No human adjudicated a single verdict
in the chain, which is the weakest point in the evidence for C1. A human-anchored
spot-check and a manual error taxonomy over stratified error instances are in progress
and are not part of this report.

\paragraph{LLM-assigned attributes.} Concept subtype, edge confidence, intervention
lifecycle and maturity are model assignments. The judge scores only node and edge
faithfulness and edge type, so the maturity profile of \cref{sec:r-corpus} is composition
rather than measurement.

\paragraph{Single extractor.} All extractions come from one model (\texttt{o3}) under one
prompt, so our design does not characterise model-specific extraction bias.
% [GAP: an n=20 multi-model consistency check (o3 / GPT-5 / Claude-4) was run earlier in
%  the project. If those numbers are recoverable, one sentence here materially
%  strengthens this limitation. Owner: Sai/Mike via ticket #150 recovery bundle.]

\paragraph{Corpus.} ARD is a curated snapshot and its coverage thins after 2023, so
frontier-era discourse is under-represented. Every claim here describes the extracted
corpus, not the field as of publication.

\subsection{Outlook}
\label{sec:outlook}

The natural next step, and the reason for the release, is cross-paper mechanism analysis:
grouping interventions into mechanism families and anchoring risks to an external risk
taxonomy, so that one can ask which mechanism families address which risks, at what stage
of validation evidence, across the whole corpus --- and, where a mechanism family is
densely populated for every risk but one, whether that gap marks a research opportunity.
We do not attempt it here, because \cref{sec:r-clustering} shows that node-level semantic
clustering on this data is too weak to carry mechanism-family grouping; doing it properly
requires an external anchor rather than a better threshold. Two further directions follow
directly from this work: the fidelity controls of \cref{sec:r-template}, and a GraphRAG
layer that answers mechanism questions from stored chains rather than from retrieved
passages.

\subsection{Conclusion}
\label{sec:conclusion}

We extracted, from 11{,}779 AI safety documents, the explicit logical chains by which
proposed interventions are argued to reduce named risks. Across 2{,}772
quality-controlled chains, 87.4\% carry all five reasoning stages while spanning 4--16
nodes, structure captured from the papers rather than imposed on them. The chains answer
mechanism-level queries that document-level resources cannot, and they expose what the
corpus holds: 2.5\% of extracted interventions are described as deployed, and one
document in six yields a complete chain at all. We also report, as a first-class result,
four measured ways in which routine graph analysis of such a corpus produces artefacts
rather than findings --- most sharply a merge-manufactured centrality hub at 90$\times$
the next node --- together with the controls that remove them, and two previously
reported findings on this corpus that do not reproduce. The chain is the unit that
survives scrutiny on this data; the aggregate graph is not.

%%%%%%%%%%%%%%%%%%%%%%%%%%%%%%%%%%%%%%%%%%%%%%%%%%%%%%%%%%%%%%%%%%%%%%%%
\section*{Acknowledgments}

This work was carried out by a distributed volunteer team convened through
\textbf{EleutherAI}, whose Discord server hosted all project coordination --- stand-ups,
working threads and review discussion --- for the full duration of the project. We thank
EleutherAI for providing and maintaining that platform.

% !! GATE -- MUST BE CLOSED BEFORE SUBMISSION. Do not remove this comment until resolved.
%
% COMPUTE ACKNOWLEDGEMENT. Model inference for the extraction and judge runs was paid for
% with API credits provided by a private acquaintance. Before this paper is submitted,
% Martin must confirm with them whether they wish to be acknowledged by name,
% acknowledged anonymously, or not mentioned at all.
%
% Why this is a gate and not a nicety: the credits may have been granted under terms that
% did not contemplate being passed to a third party. Naming the donor without asking
% could expose them to a problem with whoever issued the credits. Naming an institution
% that did not authorise the transfer would be worse.
%
% Three acceptable outcomes, pick one and delete the other two:
%   (a) named:     "Compute for this work was donated by <NAME>."
%   (b) anonymous: "Compute for this work was donated by a private benefactor who has
%                   asked not to be named."
%   (c) omitted:   no compute sentence at all (and confirm the donor is content with
%                   that -- silence is also a choice they should get to make).
%
% Default if unreachable before the deadline: (b).
%
% [GAP: compute acknowledgement -- BLOCKED on donor consent. See gate above.]

% [GAP: author list, affiliations and per-author contribution statement. Deliberately NOT
%  drafted yet -- being settled separately. Do not populate from git history alone.]

%%%%%%%%%%%%%%%%%%%%%%%%%%%%%%%%%%%%%%%%%%%%%%%%%%%%%%%%%%%%%%%%%%%%%%%%
\section*{Impact Statement}
\label{sec:impact}

This work releases a dataset and a set of measurements about a safety-research corpus.
The most likely harm from it is misreading rather than misuse: a reader who takes cluster
sizes or centrality rankings from the released graph and reports them as facts about
which risks the field neglects will produce a confident and wrong claim.
\Cref{sec:r-hub}--\cref{sec:r-bydesign} exist partly to make that harder, and the release
carries the same warning in its documentation. We make no claim about which areas of AI
safety research are under-served, and the results here do not support one.

%%%%%%%%%%%%%%%%%%%%%%%%%%%%%%%%%%%%%%%%%%%%%%%%%%%%%%%%%%%%%%%%%%%%%%%%
\section*{Use of AI Assistance}
\label{sec:aidisclosure}

Beyond the extraction and judging models, which are the object of study and are described
in \cref{sec:m-extraction,sec:m-validation}, AI tools were used in preparing this work.
An LLM coding assistant (Anthropic Claude) wrote and ran the analysis and audit scripts
released with the paper, and drafted and restructured portions of this manuscript. All
claims, all numerical results and all interpretation are the authors' own: every
quantitative statement is re-derived from the raw data by the released audit script and
was checked against it, and the authors take full responsibility for the content. No LLM
is an author, and none is credited as one.
% [GAP: scope of the drafting claim. "portions of this manuscript" is deliberately
%  imprecise until the co-authors confirm what each of them wrote by hand -- parts of
%  Methods and several appendices are reused near-verbatim from a co-author's frozen
%  Overleaf draft, and this statement must not overstate the assistant's share or
%  understate theirs. Settle alongside the author list and CRediT statement (G15); the
%  same venue policies that require this disclosure also bar an LLM from authorship, so
%  the two decisions belong in one conversation.
%  Re-check the wording against the chosen venue: ICLR 2026 desk-rejects undisclosed LLM
%  use and requires all contributions to be acknowledged; ICML 2026 permits assistance
%  but places all responsibility on the authors and forbids crediting an LLM.
%  Verified 2026-08-11.]

% Exactly ONE \bibliography command is allowed per document; the appendix must not carry
% a second one. \bibliographystyle is required -- stock article does not supply it, which
% is why citations otherwise render as "??".
\bibliographystyle{plainnat}
\bibliography{refs}

\clearpage
\appendix

\section{Extraction Prompt and Schema}
\label{app:prompt}
% [GAP: paste PROMPT_EXTRACT from
%  intervention_graph_creation/src/prompt/final_primary_prompt.py, plus the node/edge
%  Pydantic schema summary. Trim to one page.]

\section{Deduplication Threshold Experiments}
\label{app:dedup}
% Reused verbatim from Gleb_overleaf_paper.txt Appendix D.

We evaluated four threshold configurations for node deduplication, measuring the
trade-off between precision (avoiding false merges) and recall (capturing legitimate
duplicates).

\begin{table}[h]
\centering
\small
\begin{tabular}{ccccc}
\toprule
\textbf{Cosine} & \textbf{Jaccard} & \textbf{Cand. Pairs} & \textbf{Dedup \%} & \textbf{Result} \\
\midrule
0.65 & 0.30 & 416 & 0.2\% & Too conservative \\
0.85 & 0.00 & 105,390 & 11.9\% & False clusters \\
\textbf{0.88} & \textbf{0.05} & \textbf{4,411} & \textbf{1.2\%} & \textbf{Selected} \\
0.95 & 0.00 & 503 & 0.2\% & Too conservative \\
\bottomrule
\end{tabular}
\caption{Deduplication threshold experiments. Candidate pairs shows node pairs passing
the similarity threshold; after Union-Find grouping fewer nodes are removed as duplicates
merge transitively (4{,}411 pairs yield 2{,}385 removed nodes). The Jaccard-free
configuration (0.85/0.00) created false clusters merging unrelated concepts.
\Cref{sec:r-recall} reports the recall audit of the selected configuration.}
\label{tab:dedup-thresholds}
\end{table}

\paragraph{Boundary examples (selected configuration).} Representative merged pairs at
the decision boundary (cosine $\approx$ 0.88):
\begin{itemize}
\setlength{\itemsep}{0pt}
\item ``US semiconductor export controls'' $\leftrightarrow$ ``AI chip export control regulations''
\item ``Misalignment of AI objectives'' $\leftrightarrow$ ``objective misalignment in smarter-than-human AI''
\item ``Insufficient AI safety research capacity'' $\leftrightarrow$ ``Insufficient AI alignment research capacity''
\end{itemize}

\paragraph{Edge type normalisation.} The 12 canonical edge types after normalisation:
\textit{motivates, implemented\_by, validated\_by, caused\_by, mitigated\_by, enabled\_by,
addressed\_by, refined\_by, specified\_by, required\_by, preceded\_by, related\_to}.

\section{Clustering Methods Comparison}
\label{app:clustering}
% Reused verbatim from Gleb_overleaf_paper.txt Appendix F.

\begin{table}[h]
\centering
\small
\begin{tabular}{lccc}
\toprule
\textbf{Method} & \textbf{Silhouette} & \textbf{CH Score} & \textbf{DB Index} \\
\midrule
Direct 1536D, $k=5$ & 0.022 & 3,831 & 5.18 \\
Direct 1536D, $k=8$ & 0.017 & 2,922 & 5.03 \\
UMAP 150D, $k=8$ & 0.287 & 97,380 & 1.27 \\
\textbf{UMAP 150D, $k=40$} & \textbf{0.298} & \textbf{55,410} & \textbf{1.15} \\
\bottomrule
\end{tabular}
\caption{Clustering quality metrics. Direct clustering in 1536D embedding space fails
(silhouette $\approx$ 0.02). UMAP reduction to 150D enables meaningful topical clustering
with a 13--17$\times$ silhouette improvement. Calinski-Harabasz (CH) measures separation
(higher is better); Davies-Bouldin (DB) measures overlap (lower is better).}
\label{tab:clustering-methods}
\end{table}

\section{Topical Cluster Reference}
\label{app:clusters}
% Reused from Gleb_overleaf_paper.txt Appendix E. Reframed as a browsing index.
% [GAP: also publish cluster_representatives_20.json (20 representative nodes per
%  cluster) alongside the dataset so reviewers can audit cluster naming.]

The 40 topical clusters from $k$-means on UMAP-reduced embeddings, cluster sizes
2{,}065--8{,}098 nodes. These are a \emph{browsing index} over the corpus, not a taxonomy:
co-membership indicates vocabulary proximity rather than a mechanistic relationship, and
boundaries between adjacent sub-fields overlap.

\begin{small}
\begin{enumerate}
\setlength{\itemsep}{0pt}
\item[0.] RL Impact Measures, Power \& Time-Bounded Optimization
\item[1.] Foundational Alignment Research \& AGI Safety
\item[2.] Deployment Monitoring \& Safety Evaluation
\item[3.] Training Optimization \& Hyperparameter Learning
\item[4.] Information Theory \& Communication Efficiency
\item[5.] Explainability \& Interpretability Methods
\item[6.] Compute Governance \& Hardware Regulation
\item[7.] AI Ethics \& Consciousness Research
\item[8.] Compute Scaling, Recursive Self-Improvement \& Brain Preservation
\item[9.] Value Learning \& Preference Aggregation
\item[10.] Fine-Tuning Methods \& Cognitive Empathy
\item[11.] Debate \& Iterated Amplification
\item[12.] AI Governance \& Legal Frameworks
\item[13.] Data Quality \& LLM Hallucinations
\item[14.] Reward Misspecification \& Proxy Gaming
\item[15.] Curriculum Learning \& Scientific Discovery Systems
\item[16.] Economic Impact \& Technological Unemployment
\item[17.] AI Safety Community \& Education
\item[18.] Transparency \& Accountability Mechanisms
\item[19.] Corrigibility \& Shutdown Safety
\item[20.] AI Forecasting \& Timeline Prediction
\item[21.] Wellbeing-Aware Recommenders, Misuse \& Content Moderation
\item[22.] Inner \& Outer Misalignment
\item[23.] Imitation Learning \& Sample Efficiency
\item[24.] Red Teaming \& Adversarial Testing
\item[25.] Adversarial Robustness \& OOD Generalization
\item[26.] Empirical Benchmarking \& Scaling Laws
\item[27.] Reward Modeling \& Human Feedback
\item[28.] Decision Theory \& Coherence
\item[29.] Activation Engineering \& Model Steering
\item[30.] Distribution Shift \& Robustness Failures
\item[31.] Formal Verification \& Proof Systems
\item[32.] Containment \& Sandboxing Strategies
\item[33.] Biosecurity \& Critical Infrastructure
\item[34.] Legible Motion Planning, Reachability \& Safe Delegation
\item[35.] Deceptive Alignment \& Treacherous Turn
\item[36.] AI Risk Perception, Governance \& Research Prioritization
\item[37.] Funding \& Resource Allocation
\item[38.] Existential Risk \& Catastrophic Scenarios
\item[39.] Public Outreach, Framing \& Community Building
\end{enumerate}
\end{small}

\section{Similarity Threshold and Hop Sensitivity}
\label{app:sensitivity}
% Reused from Gleb_overleaf_paper.txt Appendices G + H, merged.

\paragraph{Similarity threshold ($\tau$).} We tested $\tau \in \{0.80, 0.85, 0.90,
0.95\}$ on within-category node pairs. At $\tau{=}0.80$, 169{,}083 similarity edges
connect conceptually related nodes across documents. Higher thresholds progressively
fragment the graph: $\tau{=}0.85$ reduces edges by approximately 50\%, $\tau{=}0.90$ by
over 90\%, and $\tau{=}0.95$ captures only near-duplicate phrasings. The main text uses
EDGE-only paths for all quantitative claims precisely because of this sensitivity.

\paragraph{Similarity hops.} Graph traversal with 0--3 similarity hops on 500 sampled
risks:

\begin{table}[h]
\centering
\footnotesize
\begin{tabular}{cccc}
\toprule
\textbf{Hops} & \textbf{Scope (mean)} & \textbf{Intv. (mean)} & \textbf{Intv. (med.)} \\
\midrule
0 & 14 & 2.5 & 2 \\
1 & 300 & 56 & 10 \\
2 & 6,209 & 1,172 & 120 \\
3 & 26,827 & 5,065 & 2,384 \\
\bottomrule
\end{tabular}
\caption{Similarity hop comparison. Scope is the mean number of nodes reachable from a
risk via BFS; Intv. is reachable interventions. Each hop multiplies reach by roughly
20$\times$, which is why hop-augmented reachability appears here as a sensitivity rather
than as evidence in the main text.}
\label{tab:sim-hops}
\end{table}

\paragraph{PageRank mixing coefficient ($\alpha$).} For the centrality analyses in
\cref{sec:r-hub} we combine structural and similarity transition matrices as
$P = \alpha P_{\text{struct}} + (1-\alpha) P_{\text{sim}}$. Rankings are stable across
$\alpha \in [0.5, 0.9]$ (Spearman $\rho = 1.0$, identical top-20); only pure structural
weighting ($\alpha = 1.0$) deviates, with 2 risks leaving the top-20. This stability is a
robustness property of the mixing parameter and says nothing about the merge artefact
documented in \cref{sec:r-hub}: the rankings are stably contaminated.

\section{Race-Framing Classifier Validation}
\label{app:race}
% Reused from Gleb_overleaf_paper.txt Appendix J, kept because it is a valid classifier
% validation; prevalence claims moved to the main text.

We validated keyword-based race-framing classification (keywords: ``competi-'',
``race'') on a stratified sample of 52 problem-analysis nodes. \textbf{Stratum 1}
($n{=}36$): all race-flagged PA nodes from top-100 single-path risks --- 36/36 = 100\%
precision. Example PAs: ``Competitive deployment pressures reducing safety measures,''
``AI arms race dynamics reducing safety investment.'' \textbf{Stratum 2} ($n{=}16$):
random sample from remaining risks --- 12/16 = 75\% precision. The four false positives
matched keywords but described competition \emph{among} AI systems rather than a human
race \emph{to develop} AI. Overall precision 48/52 = 92\%.

\paragraph{Recall estimation.} On a stratified sample ($n{=}50$) of non-flagged PA nodes,
one false negative appeared among top-100 single-path risks (``AI developers racing
without adequate safety''), giving recall 36/37 = 97.3\% in that stratum. Recall falls to
an estimated $\sim$68\% in tiers 101--500 due to morphological variants and synonyms
without keywords.

\paragraph{Scope.} This validates the \emph{classifier} on the sampled strata and
establishes nothing about prevalence. \Cref{sec:r-selection} reports corpus-wide
prevalence (2.21\% of problem-analysis nodes), measured separately by exhaustive scan.

\section{Judge Sample Composition}
\label{app:judge}

The 100 judged papers were sampled at random from successfully extracted papers and span
all nine ARD source types: LessWrong 28, Alignment Forum 21, blogs 15, EA Forum 11,
arXiv 11, Arbital 6, special documents 4, YouTube transcripts 2, aisafety.info 2.
% SRC: experiment_judge_item2_report.json (n=100, by_source_type) +
%      extraction_validator/STATUS.md on branch judge_handoff_workshop_items_2_3

The composition matters for reading \cref{sec:r-judge}: the audited sample is
forum-weighted in the same way the corpus is, so the error profile it yields describes
extraction over heterogeneous web documents rather than over arXiv papers alone.

\section{A Failure Case}
\label{app:failure}

Quality cuts do not make every extracted chain meaningful. The chain below passes every
structural filter we apply --- seven nodes, all five reasoning stages, edge confidence
$\geq 3$, a maturity-4 intervention endpoint, single source --- and is nonetheless
spurious. Its source is an expository mathematics article proving that there are
infinitely many primes.

\begin{quote}\small
\textbf{risk} Prime scarcity affecting cryptographic key space $\rightarrow$
\textbf{pa} Finite prime assumption in key generation algorithms $\rightarrow$
\textbf{ti} Contradiction from finite prime list assumption $\rightarrow$
\textbf{dr} Product-plus-one construction to find new prime $\rightarrow$
\textbf{im} Compute $P = \prod(\text{all known primes}) + 1$ $\rightarrow$
\textbf{va} Proof that $P$ has a prime divisor outside the original list
$\rightarrow$ \textbf{intv} Integrate infinite-prime theorem into cryptographic key
generation standards
\end{quote}
% SRC: node ids [60117..60123], source https://arbital.com/p/infinitely_many_primes

The extraction has taken Euclid's proof and cast it into the risk--intervention schema,
inventing a safety framing (``prime scarcity'') that the source never asserts. This is
the characteristic failure mode of schema-guided extraction over a heterogeneous corpus:
when a document contains a problem, an argument and a construction, the schema will find
a chain whether or not the document concerns safety. Note which failure this is. The
source is a real mathematical argument, so the failure is not that the pipeline read an
out-of-domain document, but that it \emph{invented} a risk framing absent from the
source. That is a fidelity failure of the kind \cref{sec:r-template} identifies as
needing a control, and it is why the structural statistics there are reported as
structural properties and never as evidence that a chain is substantively correct.
Establishing the rate at which this happens requires the schema-ablation and
degraded-source controls; this appendix is an existence proof only.

% [GAP: optionally add 1-2 further chains (e.g. one with a missing stage) if space
%  permits. The failure case above is the one that must stay.]

% NOTE: no \bibliography here -- it lives once, before the appendix. A second one makes
% BibTeX abort with "Illegal, another \bibdata command".

\end{document}
